\documentclass[11pt]{article}
\usepackage{amsmath,amssymb,amsthm,latexsym}
\usepackage{mathrsfs,mathtools}
\usepackage{graphicx}
\usepackage{comment}
\usepackage{bm}
\usepackage{natbib}
\usepackage{subfigure,multirow}
\newtheorem{theorem}{Theorem}
\newtheorem{lemma}{Lemma}
\usepackage[colorlinks,linkcolor=black,citecolor=black,urlcolor=black]{hyperref}
\usepackage{verbatim}

\renewcommand\baselinestretch{1.25}
\topmargin -.5in
\textheight 9in
\textwidth 6.6in
\oddsidemargin -.25in
\evensidemargin -.25in

\title{Solutions to Homework 7}
\author{Xiaochuan Gong}
\date\today

\begin{document}
\maketitle

\subsection*{Problem 1}
Assume $\nabla f\in\text{Lip}(L)$ and $f$ is strongly convex with parameter $m$, prove that gradient descent with step size $\delta_k=\frac{2}{L+m}$ satisfies 
\[
f(x_k)-f(z)+\frac{m}{2}||x_k-z||^2 \leq \left(\frac{L-m}{L+m}\right)^k\left(f(x_0)-f(z)+\frac{m}{2}||x_0-z||^2\right),\ \forall z.
\]

\begin{proof}
First, since $\nabla f\in\text{Lip}(L)$, we have 
\begin{equation}
    \label{1}
    f(y) \leq f(x)+\langle\nabla f(x),y-x\rangle+\frac{L}{2}||y-x||^2,\ \forall x,y\in\mathbb{R}^d. 
\end{equation}
We let $y=x-\frac{2}{L+m}\nabla f(x)$ and plug $y$ in \eqref{1}, we get
\begin{equation}
    \label{2}
    f(y) \leq f(x)-\frac{2m}{(L+m)^2}||\nabla f(x)||^2.
\end{equation}
Also recall that if $f$ is strongly convex with parameter $m$, there exists
\begin{equation}
    \label{3}
    f(y) \geq f(x)+\langle\nabla f(x),y-x\rangle+\frac{m}{2}||y-x||^2,\ \forall x,y\in\mathbb{R}^d. 
\end{equation}
For any $z\in\mathbb{R}^d$, we have
\begin{equation*}
    \begin{aligned}
        \frac{m}{2}||y-z||^2
        & = \frac{m}{2}||x-\frac{2}{L+m}\nabla f(x)-z||^2 \\
        & = \frac{m}{2}||x-z||^2 - \frac{2m}{L+m}\langle\nabla f(x),x-z\rangle + \frac{2m}{(L+m)^2}||\nabla f(x)||^2 \\
        & \leq \frac{m}{2}||x-z||^2 + \frac{2m}{m+L}\left[f(z)-f(x)-\frac{m}{2}||z-x||^2\right] + \frac{2m}{(L+m)^2}\cdot\frac{(L+m)^2}{2m}(f(x)-f(y)) \\
        & = \frac{m(L-m)}{2(L+m)}||x-z||^2 - \frac{2}{L+m}(f(x)-f(z)) + f(x)-f(y), \\
    \end{aligned}   
\end{equation*}
where  we use \eqref{2} and \eqref{3} to derive the inequality step. Then we add $f(y)-f(z)$ to both sides of the above inequality, we obtain
\[
f(y)-f(z)+\frac{m}{2}||y-z||^2 \leq \left(\frac{L-m}{L+m}\right)\left(f(x)-f(z)+\frac{m}{2}||x-z||^2\right),\ \forall z.
\]
Set $x=x_{i-1}$, $y=x_i$ for $i=1,\dots,k$ and by recursion we get
\[
f(x_k)-f(z)+\frac{m}{2}||x_k-z||^2 \leq \left(\frac{L-m}{L+m}\right)^k\left(f(x_0)-f(z)+\frac{m}{2}||x_0-z||^2\right),\ \forall z.
\]
\end{proof}

\subsection*{Problem 2}
Assume $\nabla f\in\text{Lip}(L)$ and $f$ is strongly convex with parameter $m$. Prove the following convergence rates. \\
\noindent (a)\ For gradient flow $\frac{d}{dt}X_t=-\nabla f(X_t)$, let $\mathcal{E}_t=e^{\frac{2mL}{m+L}}\frac{1}{2}||x^*-X_t||^2$ be the Lyapunov function, then we have $f(X_t)-f(x^*)\leq O(\frac{L}{2}e^{-\frac{2mL}{m+L}t})$. \\
\noindent (b)\ For gradient descent $\frac{x_{k+1}-x_k}{\delta}=-\nabla f(x_k)$ with $0<\delta\leq\frac{2}{m+L}$, let $E_k=\left(1-\frac{2mL}{m+L}\delta\right)^{-k}\frac{1}{2}||x^*-x_k||^2$ be the Lyapunov function, then we have $f(x_k)-f(x^*)\leq O(\frac{L}{2}e^{-\frac{2mL}{m+L}\delta k})$.

\begin{proof}
(a)\ First, $\nabla f\in\text{Lip}(L)$ and $f$ is $m$-strongly convex ($m\leq L$), we have
\begin{equation}
    \label{4}
    \langle \nabla f(X_t),x^*-X_t \rangle \leq -\frac{mL}{m+L}||x^*-X_t||^2-\frac{1}{m+L}||\nabla f(X_t)||^2.
\end{equation}
We provide a proof of this bound in the Appendix \ref{A}. Using \eqref{4}, it follows that 
\begin{equation*}
    \begin{aligned}
        \frac{d}{dt}\mathcal{E}_t 
        & = e^{\frac{2mL}{m+L}t}\left(\frac{mL}{m+L}||x^*-X_t||^2-\langle\frac{d}{dt}X_t,x^*-X_t\rangle\right) \\
        & = e^{\frac{2mL}{m+L}t}\left(\frac{mL}{m+L}||x^*-X_t||^2+\langle\nabla f(X_t),x^*-X_t\rangle\right) \\
        & = e^{\frac{2mL}{m+L}t}\left(\frac{mL}{m+L}||x^*-X_t||^2-\frac{mL}{m+L}||x^*-X_t||^2-\frac{1}{m+L}||\nabla f(X_t)||^2\right) \\
        & \leq 0.
    \end{aligned}
\end{equation*}
By integrating, we obtain the statement 
\[
\mathcal{E}_t-\mathcal{E}_0 = \int_{0}^{t}\frac{d}{dt}\mathcal{E}_sds \leq 0.
\]
Thus we have
\[
\frac{1}{2}||x^*-X_t||^2 \leq e^{-\frac{2mL}{m+L}t}\mathcal{E}_0.
\]
In addition, by the property of $L$-smoothness \eqref{1}, we can subsequently conclude the upper bound for the optimality gap as follows:
\begin{equation*}
    \begin{aligned}
        f(X_t)-f(x^*) & \leq \langle\nabla f(x^*),X_t-x^*\rangle+\frac{1}{2}L||x^*-X_t||^2 \\
                      & = \frac{1}{2}L||x^*-X_t||^2 \\
                      & \leq e^{-\frac{2mL}{m+L}t}\frac{L}{2}||x^*-X_0||^2.
    \end{aligned}
\end{equation*}
Hence we obtain $f(X_t)-f(x^*)\leq O(\frac{L}{2}e^{-\frac{2mL}{m+L}t})$.
\\

\noindent (b)\ For gradient descent, as long as the function is $L$-smooth and $m$-strongly convex, where $\delta\leq\frac{2}{m+L}$, the following function, 
\[
E_k = \left(1-\frac{2mL}{(m+L)}\delta\right)^{-k}\frac{1}{2}||x^*-x_k||^2,
\]
is a Lyapunov function. We check,
\begin{equation*}
    \begin{aligned}
        \frac{E_{k+1}-E_k}{\delta}
        & = \left(1-\frac{2mL}{(m+L)}\delta\right)^{-(k+1)}\frac{1}{\delta}\left[\frac{1}{2}||x^*-x_{k+1}||^2-\frac{1}{2}\left(1-\frac{2mL}{m+L}\delta\right)||x^*-x_k||^2\right] \\
        & = \left(1-\frac{2mL}{(m+L)}\delta\right)^{-(k+1)}\left[\frac{\left(\frac{1}{2}||x^*-x_{k+1}||^2-\frac{1}{2}||x^*-x_{k}||^2\right)}{\delta}+\frac{mL}{m+L}||x^*-x_k||^2\right] \\
        & = \left(1-\frac{2mL}{(m+L)}\delta\right)^{-(k+1)}\left(-\langle\frac{x_{k+1}-x_k}{\delta},x^*-x_k\rangle+\varepsilon_k^1+\frac{mL}{m+L}||x^*-x_k||^2\right) \\
        & = \left(1-\frac{2mL}{(m+L)}\delta\right)^{-(k+1)}\left(\langle\nabla f(x_k),x^*-x_k\rangle+\varepsilon_k^1+\frac{mL}{m+L}||x^*-x_k||^2\right) \\
        & \leq \left(1-\frac{2mL}{(m+L)}\delta\right)^{-(k+1)}\left(-\frac{mL}{m+L}||x^*-x_k||^2+\varepsilon_k^2+\frac{mL}{m+L}||x^*-x_k||^2\right) \\
        & = \left(1-\frac{2mL}{(m+L)}\delta\right)^{-(k+1)}\varepsilon_k^2 \\
        & \leq 0,
    \end{aligned}
\end{equation*}
where $\varepsilon_k^1=\frac{\delta}{2}||\nabla f(x_k)||^2$ and $\varepsilon_k^2=-\left(\frac{1}{m+L}-\frac{\delta}{2}\right)||\nabla f(x_k)||^2$. We also use \eqref{4} to derive the first inequality. Since we take $\delta\in\left(0,\frac{2}{m+L}\right]$, then $\varepsilon_k^2\leq0$. By summing, we obtain the statement $E_k-E_0\leq\sum_{i=0}^{k}\frac{E_{i+1}-E_i}{\delta}\delta\leq0$, thus we have
\[
\frac{1}{2}||x^*-x_k||^2 \leq \left(1-\frac{2mL}{(m+L)}\delta\right)^{k}E_0.
\]
In addition, by the property of $L$-smoothness \eqref{1}, we can subsequently conclude the upper bound for the optimality gap as follows:
\begin{equation*}
    \begin{aligned}
        f(x_k)-f(x^*) & \leq \langle\nabla f(x^*),x_k-x^*\rangle+\frac{1}{2}L||x^*-x_k||^2 \\
                      & = \frac{1}{2}L||x^*-x_k||^2 \\
                      & \leq e^{-\frac{2mL}{m+L}\delta k}\frac{L}{2}||x^*-x_0||^2.
    \end{aligned}
\end{equation*}
Hence we obtain $f(x_k)-f(x^*)\leq O(\frac{L}{2}e^{-\frac{2mL}{m+L}\delta k})$. 
\end{proof}


\subsection*{Problem 3}
Consider the proximal operators of matrices
\[
\text{Prox}_h(Y) \coloneqq \arg\min_{X\in\mathbb{R}^{m\times n}}\left\{\frac{1}{2}||X-Y||_{F}^2+h(X)\right\}. 
\]
Let $h(X)=\lambda||X||_*$. Show that the proximal operator of the nuclear norm is 
\[
\text{Prox}_h(Y) = U\text{diag}(\{(\sigma_i-\lambda)_{+}\}_{1\leq i\leq r})V^{T},
\]
where $Y=U\Sigma V^{T},\ \Sigma=\text{diag}(\{\sigma_i\}_{1\leq i\leq r})$ is the SVD of $Y\in\mathbb{R}^{m\times n}$ of rank $r$, $\lambda\geq0$, $t_{+}=\max(0,t)$, $||\cdot||_F$ denotes the Frobenius norm and $||\cdot||_*$ denotes the nuclear norm. 

\begin{proof}
Since $h_0(X)\coloneqq\frac{1}{2}||X-Y||_{F}^2+\lambda||X||_*$ is strictly convex, it easy to see that there exists a unique minimizer, and we thus need to prove that it is equal to $U\text{diag}(\{(\sigma_i-\lambda)_{+}\}_{1\leq i\leq r})V^{T}$. We denote $\hat X=\text{Prox}_h(Y)$ for convenience. Note that $\hat X$ minimizes $h_0(X)$ if and only if $0\in\partial h_0(X)$, i.e.,
\begin{equation}
    \label{6}
    0 \in \hat X-Y+\lambda\partial||\hat X||_*, 
\end{equation}
where $\partial||\hat X||_*$ is the subdifferential of the nuclear norm. Let $X\in\mathbb{R}^{m\times n}$ be an arbitary matrix and $U\Sigma V^{T}$ be its SVD. Recall that the subdifferential of the nuclear norm is
\[
\partial||X||_* = \left\{UV^{T} + W\ | \ U^{T}W = 0, WV = 0, ||W||_2 \leq 1, W\in\mathbb{R}^{m\times n}\right\},
\]
where $||\cdot||_2$ denotes the spectral norm. Now we set $\hat X\coloneqq U\text{diag}(\{(\sigma_i-\lambda)_{+}\}_{1\leq i\leq r})V^{T}$ for short. In order to show that $\hat X$ satisfies \eqref{6}, we decompose the SVD of $Y$ as 
\[
Y = U_1\Sigma_1V_1^{T}+U_2\Sigma_2V_2^{T},
\]
where $U_1$, $V_1$ are the singular vectors associated with singular values greater than $\lambda$, and $U_1$, $V_1$ are the singular vectors associated with singular values smaller or equal to $\lambda$. With these notations, we have 
\[
\hat X = U_1(\Sigma_1-\lambda I)V_1^{T},
\]
and, therefore, 
\[
Y-\hat X = \lambda(U_1V_1^{T}+W), \quad W = U_2(\frac{1}{\lambda}\Sigma_2)V_2^{T}.
\]
By definition, the columns of $U$ and $V$ are orthonormal. Thus we have $U_1^{T}W=\frac{1}{\lambda}U_1^{T}U_2\Sigma_2V_2^{T}=0$, and $WV_1=\frac{1}{\lambda}U_2\Sigma_2V_2^{T}V_1=0$. Since the diagonal elements of $\frac{1}{\lambda}\Sigma_2$ have magnitudes bounded by 1, we also have $||W|||_2\leq1$. Hence $Y-\hat X \in \lambda\partial||\hat X||_*$, which concludes the proof. 
\end{proof}


\section{Appendix}\label{A}
We provide a proof of inequality \eqref{4} there.

\begin{proof}
Define $\phi(x)=f(x)-\frac{m}{2}||x||^2$ and note $\nabla\phi(x)=\nabla f(x)-mx$. Since $f$ if $L$-smooth, then $\langle\nabla f(x)-\nabla f(y),x-y\rangle \leq L||x-y||^2$, and this implies $\langle\nabla\phi(x)-\nabla\phi(y),x-y\rangle \leq (L-m)||x-y||^2$. Thus $\phi$ is $(L-m)$-smooth. This, in turn, implies
\[
\langle\nabla\phi(x)-\nabla\phi(y),x-y\rangle \geq \frac{1}{L-m}||\nabla\phi(x)-\nabla\phi(y)||^2.
\]
Expanding the above inequality, and we obtain
\[
\langle\nabla f(x)-\nabla f(y),x-y\rangle \geq \frac{mL}{m+L}||x-y||^2+\frac{1}{m+L}||\nabla f(x)-\nabla f(y)||^2.
\]
Let $x=X_t$, $y=x^*$, and note $\nabla f(x^*)=0$, then we get
\[
\langle\nabla f(X_t),X_t-x^*\rangle \geq \frac{mL}{m+L}||X_t-x^*||^2+\frac{1}{m+L}||\nabla f(X_t)||^2.
\]
Multiply both sides by -1 and we obtain the inequality \eqref{4}.
\end{proof}

\end{document}